\documentclass[11pt,letterpaper]{article}
\usepackage[T1]{fontenc}
\usepackage{lmodern,amsmath,amssymb,amsthm,microtype}
\usepackage[margin=1in]{geometry}
\usepackage[hidelinks]{hyperref}
\providecommand{\tightlist}{\setlength{\itemsep}{0pt}\setlength{\parskip}{0pt}}
\DeclareUnicodeCharacter{2264}{\ensuremath{\le}}
\DeclareUnicodeCharacter{2265}{\ensuremath{\ge}}
\DeclareUnicodeCharacter{2260}{\ensuremath{\ne}}
\DeclareUnicodeCharacter{2212}{\ensuremath{-}}

\DeclareUnicodeCharacter{2211}{\ensuremath{\sum}}

\usepackage{longtable,booktabs,array}
\title{A 104-shard exclusion for the penultimate real subset-sum weight}
\author{Aafi Khan Malek\\ \small Independent researcher, Ahmedabad, India\\
\small\href{mailto:aafimalek2032@gmail.com}{aafimalek2032@gmail.com}}
\date{28 November 2026}
\begin{document}
\maketitle
\textbf{Theorem.} Let nine sorted positive real weights have pairwise
unit-separated subset sums. If \(149\le x_9<150\), then \(x_8\ge134\).
The proof uses the preceding
\href{https://github.com/Aafimalek/novel-discoveries/tree/7cfbde384eb7649281ab0d31c19394540c504940/nine-separated-subset-sums/verification/rank-cuts-below-150/penultimate-cap-132}{anchored-tail exclusion}
\(x_8\ge133\) and exact finite arithmetic with no floating-point tests.
This does not determine the optimum of the nine-weight problem.

\section{Cap-specific implications}\label{cap-specific-implications}

For cap \(t=133\), the eight Harper runs force starts 144 through 172,
so the full-sign slots are \(q=36,\dots,64\). The magnitudes
\(\alpha_q=u_{a-q+1}\in[149-2q,151-2q)\) are consecutive, with constant
\(64\le a\le74\) and \(b=148-a\). The capacity ranges are:

\begin{longtable}[]{@{}lr@{}}
\toprule\noalign{}
Half-open magnitude range & Capacity \\
\midrule\noalign{}
\endhead
\bottomrule\noalign{}
\endlastfoot
{[}1,21) & \(a-64\) \\
each {[}149−2q,151−2q), q from 64 down to 36 & 1 \\
{[}79,221) & 35 \\
{[}278,321) & \(148-2a\) \\
{[}377,10000) & \(a-20\) \\
\end{longtable}

The two wide gaps each have width at least
\(\alpha_{36}-\alpha_{64}+4\ge60\). Thus signs may be flipped on weights
strictly below 30 or swapped at adjacent gaps strictly below 30 while
preserving membership in the rank groups at the two cuts. The small-flip
orbit count forces \(x_3\ge30\) because the number \(s\) of weights
below 30 has \(2^s\mid a,b\) and \(a+b=148\). A nonsmall coordinate
exists since otherwise all magnitudes would be below \(8\cdot30<377\).
The same orbit-size and interval-conditioned orbit assignment rules as
in the preceding cap packet apply, including an empty middle group when
\(a=b=74\).

For \(y=(127,\dots,134)\) and deficits \(r_i=y_i-x_i\), let
\(R=\sum r_i\). A prescribed unbalanced magnitude needs
\(R>246-79=167\), because unbalanced signs at y have magnitude at least
246. If all prescribed signs are balanced, their cap magnitude is at
most 16, so \(\alpha_{36}\ge77\) requires \(R\ge61\). Thus \(R\ge61\) in
either case. The prior cap-132 theorem supplies \(x_8\ge133\), hence
\textbf{\(r_8\le1\)}. As in the predecessor proof, subtract \(r_8\) from
each ordered tail deficit and bound the residual signed tail with prefix
averages. Adding the common contribution back gives
\(\mu T + (\sum\sigma-N\mu)r_8\) for a tail of length N and total T. The
exact bounds use \(0\le r_8\le\min(1,T_{\rm upper}/N)\) and round signed
contributions outward.

\section{Deterministic checkpoint and completed
arithmetic}\label{deterministic-checkpoint-and-completed-arithmetic}

Of 8,001 initial unit cells, 7,095 reached the count splits. The 78,045
count profiles lose 37,660 to pure orbit sizes, 4,862 to ordinary
interval tests, and 23,011 to the interval-conditioned orbit assignment.
The exact phase-2 replay generated \textbf{12,512} residual profiles,
split into \textbf{104 disjoint exhaustive shards} of at most 128
profiles by the original count and list position.

All \textbf{104} clear by exact refinement: 81 at depth 3 and 23 at
depth 4. Across the guarded worker logs:

\begin{center}
\small
\begin{tabular}{rrrrrr}
\toprule
Depth & Cells & Orbit size & Intervals & Orbit + intervals & Surviving \\
\midrule
3 & 590,986 & 158,311 & 88,006 & 343,247 & 1,422 \\
4 & 45,865 & 35,642 & 4,173 & 6,050 & 0 \\
\bottomrule
\end{tabular}
\end{center}

An independent manifest audit checks the exact domain partition, initial
payload hash, source pin and each residual SHA chain. A separate
lightweight count check reconciles all 127 recorded depths against their
exact arithmetic logs. The hashes preserve the exact artifact and do not
themselves prove the inequalities.

\section{Proof replay}\label{proof-replay}

\begin{verbatim}
sha256sum -c SHA256SUMS
python3 coverage_audit.py
python3 arithmetic_totals.py
./replay-all.sh
\end{verbatim}

The final script copies the source and frozen phase-2 checkpoint to a
fresh temporary directory. It independently regenerates phase 2, demands
a byte-identical canonical payload, and recomputes all 104 shards
serially into fresh progress. Each invocation is guarded by 35 seconds
and 143,000 KiB of virtual memory. Only one arithmetic worker runs at a
time. If disk or memory reserves are exhausted, the shell stops rather
than claiming results for missing shards. All arithmetic uses the Python
standard library with exact integers and fractions; the geometric
implications are handwritten and no Lean certification of this cap is
claimed.

\end{document}
